% Generated draft: requires mathematical and source review.
\documentclass[11pt]{article}
\usepackage[margin=25mm]{geometry}
\usepackage{fontspec}
\setmainfont{Latin Modern Roman}
\usepackage{amsmath,amssymb,mathtools,unicode-math}
\setmathfont{Latin Modern Math}
\usepackage{enumitem,needspace,xurl,hyperref}
\setlength{\parindent}{0pt}
\setlength{\parskip}{5pt}
\setlength{\emergencystretch}{3em}
\newcommand{\N}{\mathbb{N}}
\newcommand{\Z}{\mathbb{Z}}
\newcommand{\Q}{\mathbb{Q}}
\newcommand{\R}{\mathbb{R}}
\newcommand{\C}{\mathbb{C}}
\newcommand{\F}{\mathbb{F}}
\newcommand{\A}{\mathbb{A}}
\newcommand{\PP}{\mathbb{P}}
\newcommand{\E}{\mathbb{E}}
\newcommand{\eps}{\varepsilon}
\newcommand{\dd}{\,\mathrm{d}}
\begin{document}
\section*{UnsolvedMath ID: 30004820}
\textbf{Source problem code:} catalogue entry
\subsection{Definitions and mathematical statement}
Write \(\N=\{1,2,\ldots\}\). The monoid \(\operatorname{Inc}(\N)\) consists of strictly increasing maps \(\pi:\N\to\N\), with composition as its operation. For a standard graded polynomial ring \(S\) over a field \(K\) and a nonzero finitely generated graded \(S\)-module \(M\), its Castelnuovo--Mumford regularity is
\[
\operatorname{reg}_S(M)=\max\{j-i:\operatorname{Tor}_i^S(M,K)_j\ne0\}.
\]
In the question below, regularity is taken for the ideal as a graded module, as indicated by the notation \(\operatorname{reg}(I_n)\).

The supplied record does not specify the ambient graded rings \(R_n\), the coefficient field, the action of \(\operatorname{Inc}(\N)\), the precise filtration and stability conventions, or the definition of the weight \(w(I_r)\). These are unresolved parts of the formulation; the notation below retains their source-dependent meanings and does not supply missing hypotheses.

The recorded question is the following. Let \(I=(I_n)_{n\ge1}\) be an \(\operatorname{Inc}(\N)\)-invariant filtration with stability index \(r\), where each \(I_n\subset R_n\) is a nonzero proper homogeneous ideal. Must there exist integers \(a,b\) and \(N\ge r\), depending on the filtration, such that
\[
\operatorname{reg}_{R_n}(I_n)=an+b
\qquad\text{for every integer }n\ge N?
\]
For every such filtration consisting of monomial ideals, must there exist a constant \(D\), independent of \(n\), and an integer \(N\ge r\) such that
\[
\operatorname{reg}_{R_n}(I_n)=\bigl(w(I_r)-1\bigr)n+D
\qquad\text{for every integer }n\ge N?
\]
The prescribed slope in the second question cannot be interpreted numerically until \(w\) is defined. The supplied record specifically reports that eventual linearity remains open even for nonsaturated filtrations of squarefree monomial ideals. Here a squarefree monomial ideal means an ideal generated by monomials whose variable exponents belong to \(\{0,1\}\); the meaning of saturation of the filtration is not supplied. This reported open status is a source qualification, not an independently verified conclusion.
\subsection{Short English statement}
The question asks whether regularity is eventually an affine function of the index for every specified equivariant filtration, and whether monomial filtrations have the particular slope \(w(I_r)-1\). The supplied record includes nonsaturated squarefree monomial filtrations among the reported open cases. Missing ambient-ring, action, stability, saturation, and weight definitions require human review before this can be a fully specified standalone statement.
\subsection{Sources}
\begin{itemize}
\item Supplied local UnsolvedMath registry record, permanent ID \(30004820\), source problem code \texttt{catalogue entry}, titled \emph{Regularity of Equivariant Monomial-Ideal Filtrations}.
\item Supplied source citation: \emph{MFO-RIMS Tandem Workshop: Symmetries on Polynomial Ideals and Varieties (2022)}, Oberwolfach Reports, pp.~2365--2396, DOI \texttt{10.4171/owr/2021/43}. Supplied source and literature link: \url{https://doi.org/10.4171/owr/2021/43}.
\item Supplied attribution: Hop D. Nguyen, \emph{Projective dimension and regularity of equivariant filtrations of ideals}, joint work with Dinh Van Le. The record lists proposed year \(2021\), but does not identify a proposer separately.
\item Supplied literature reference labeled Le--Nguyen: \url{https://arxiv.org/abs/2206.15141}.
\item Literature checked at: \texttt{2026-08-22}. This is supplied dated curation, not independently verified. The supplied assessment labels the problem open, says the cited authors leave eventual linearity open even for nonsaturated squarefree monomial filtrations, and describes later results as covering only selected saturated or asymptotic cases. The background explicitly requests independent verification of time-sensitive status and cited claims.
\item The registry labels the statement \texttt{corrected\_verified} and reports that the official report was checked directly. These are supplied verification claims; no source consultation or independent verification is asserted here.
\end{itemize}
\end{document}
